非平滑接触动力学（NSC）通过采用速度级的时间步进法（time-stepping schemes）与线性互补问题（Linear Complementarity Problem, LCP），成为计算机图形学、机器人仿真及工程多体动力学中处理刚性碰撞的核心架构之一\cite{pfeiffer2012nonsmooth,stewart1996implicit,anitescu1997formulating}。传统的接触检测通常依赖于多边形网格（polygon mesh），并通过 GJK 等边界距离算法计算穿透深度与极值法向\cite{gilbert1988fast}。

近年来，随着几何计算与稀疏体素存储的发展，利用有向距离场（Signed Distance Field, SDF）作为刚体表面的隐式代理（implicit surrogate）成为了新的热点\cite{frisken2000adf,museth2013vdb,museth2021nanovdb,macklin2020sdfcollision}。然而，传统动力学引擎若只提取 SDF 值 $\Phi$ 与局部一阶梯度 $\nabla \Phi$，仍然面临着不可回避的数学退化问题。具体表现为：当物体接触表面存在较高曲率（如凸轮尖端或齿轮渐开线边缘）时，网格体素离散化与亚体素插值误差会被时间积分放大为畸变的接触冲量，从宏观上表现为速度求解中的高频振荡、毛刺与微弹跳。

若试图直接依靠成倍提高空间分辨率（即减小体素尺寸 Voxel Size）来平抑误差，不仅存在 $\mathcal{O}(N^3)$ 的内存暴涨灾难，且从本质上无法解决大时间步长下沿一阶切平面过度滑动所诱发的隐式侵入（Penetration）失真。

针对上述挑战，本文不再假设单一接触代理能够统一处理所有刚体接触，而是将高阶 SDF 接触建模建立在“接触类型感知”的思路之上。对于光滑连续滑移接触，本文采用带有 persistent manifold 的 \texttt{SlidingPatch} 表示，并在该局部流形上注入二阶曲率前馈；对于紧凑主导点或非光滑奇异边界接触，本文则采用更保守的 \texttt{CompactDiscrete} 表示，以抑制由分支切换和尖锐拓扑带来的高阶不稳定性。其核心仍然是利用距离场的二阶张量（Hessian）度量表面的局部曲率，并将相应的时间积分域预测量以前馈位移偏置的方式注入 Jean--Stewart--Anitescu 一类速度级互补约束中\cite{jean1999nonsmooth,stewart1996implicit,anitescu1997formulating,anitescu2004constraint}。

本工作的主要贡献包括：
\begin{itemize}
    \item \textbf{接触类型感知的高阶 SDF 接触框架}：提出一种按接触类型选择接触表示与高阶几何修正的 SDF 接触建模框架。本文不再将单一接触代理硬性施加于所有场景，而是将接触问题区分为光滑连续滑移接触与紧凑/奇异主导接触两类，并分别采用 \texttt{SlidingPatch} 与 \texttt{CompactDiscrete} 两种局部流形表示，使高阶几何信息的使用建立在与接触类型相匹配的接触表征之上。
    \item \textbf{面向连续接触的 persistent sliding manifold 与二阶前馈建模}：在光滑连续接触场景下，构建了结合 representative-query、局部曲率前馈、面向 onset 的候选筛选与局部几何拟合的持续滑移流形建模方法，使二阶曲率补偿不仅作用于稳态滑移阶段，也能覆盖首次接触建立到持续贴合保持的全过程，从而提升连续接触中的轨迹贴合能力与数值稳定性。
    \item \textbf{面向非光滑接触的紧凑离散 patch 表示与边界分析}：在紧凑主导点及非光滑奇异边界接触场景下，提出与之对应的 \texttt{CompactDiscrete} 接触策略，通过 patch 聚类、法向一致性约束与主导几何保留抑制复杂啮合中的非物理脉冲；同时结合 Twin Gear 算例揭示了二阶点接触模型在 $C^0$ 奇异边界附近仍然存在的结构性误差下限。
    \item \textbf{覆盖解析、工业与应力场景的系统性验证}：建立了由工业凸轮、偏心圆盘-滚子解析基准、Onset Stress Benchmark 以及 Twin Gear 啮合算例组成的验证体系。实验结果表明：高阶 SDF 方法在光滑连续接触中能够取得明确精度收益，在接触建立与持续滑移问题中表现出更好的稳定性，而在非光滑奇异接触中虽仍能优于一阶模型，但其性能受限于局部点接触微分近似的适用边界。
\end{itemize}
